\section{Discussion}
\label{sec:discussion}

\paragraph{What the prespecified tests show}
Under a margin fixed before the curve was computed, four of the 16 GRABMyo electrodes could be removed while remaining non-inferior to the full array. Five could not. Every $k$ from 12 to 15 met the margin, so the result does not rest on a single lucky crossing. The margin is an engineering tolerance, not a clinical threshold. Every reduction cost some Macro-F1, and halving the array cost about \DeltaKEight. The practical reading is modest: in a subject-independent setting, this 16-electrode array contains only limited redundancy for 16-gesture recognition. User-dependent studies that retain far fewer channels answer a different question~\cite{phinyomark2018emg,qu2021reduce}, and their results should not be taken as evidence that a population-level layout can be equally sparse.

Two of the three hypotheses about attribution were supported. Attribution-ranked subsets beat the average random subset at every size tested (H2), and deletion with retraining showed the expected ordering at every size (H4). The ranking is therefore not decorative: it reflects what the retrained models depend on. The stability penalty, however, did nothing (H3). With cross-fitting over roughly 34 participants per inner-training set, inner-fold attributions already agreed so closely that the dispersion term was a few percent of the magnitude term. It could not reorder channels in a way that mattered. We report SR-CFA as a tested and unsupported variant. The supported method is plain cross-fitted TreeSHAP magnitude, and a related robustness-weighted score has been proposed independently~\cite{xiang2026roshap}.

\paragraph{Attribution is not better than simple alternatives}
The exploratory comparisons temper H2. Beating random subsets is a low bar. A data-free geometric rule that spreads electrodes around both rings matched the attribution ranking at 4, 8 and 12 electrodes, and permutation importance matched it at three of four sizes. Attribution had a clear advantage only at $k = 6$. Within each participant's own random subsets, the attribution subset was typically above the median but rarely near the top. This fits a picture in which even circumferential coverage explains most of what any ranking can achieve, and the choice among well-spread layouts matters little. For a device designer, the evidence favours a well-spread layout of about 12 electrodes over trusting any particular ranking to identify an unusually good sparse subset.

\paragraph{Interpretation of the ranked electrodes}
The four highest-ranked electrodes (\TopFourSet) were the same in \TopFourFolds\ of 43 folds and in single-participant rankings (mean Spearman $\rho$ = \StabSubjRho). We deliberately do not map them to named muscles. The dataset documents ring positions, not the anatomical placement of each circumferential index. Surface electrodes also record compound activity from several muscles at each site~\cite{farina2004extraction}. The robust statement is positional. High attribution concentrated at a few circumferential positions (Fig.~\ref{fig:attribution}a). The top and bottom of the ranking were stable, and the ordering varied mainly among mid-ranked electrodes. This fits the finding that many different well-spread subsets perform similarly.

\paragraph{Cross-session and replication}
Cross-session results were exploratory and should be read narrowly. Held-out participants were not recognised worse on later days than on the training day. This shows only that no extra cross-day loss was detectable on top of the cross-user gap, not that electrode shift is harmless~\cite{hargrove2006effect,kaufmann2010fluctuating,vidovic2016improving}. The gain from training on two sessions is confounded with doubling the training data. On DB5 the whole procedure ran without modification, but the replication was underpowered. Ten subjects and a low baseline left the $k = 12$ decision on the margin, with neighbouring $k$ failing. It is best read as a non-contradiction rather than a replication of the reduction.

\paragraph{Limitations}
The primary unit is a monopolar electrode. A conventional device measures bipolar pairs, and removing four electrodes does not remove four sensors of such a device. The exploratory bipolar re-analysis points the other way: with pairs as the unit, not even one of eight could be removed within the primary margin, although the shortfall at seven pairs was small. The 25\% reduction should therefore be read as a statement about monopolar electrodes on this array, not as evidence that a bipolar armband can shed a quarter of its sensors. The bipolar reference is itself a derived montage, and physical bipolar electrodes may behave differently. All participants were able-bodied, and all analysis was offline. The SVM was trained on a subsample and is a lower bound. The H2 and H4 comparisons used a lighter comparison model than the headline curve. The count-feature threshold was a single scalar set from pooled data. $\lambda$ was not selected, and the detailed H2--H4 test families were specified after the underlying results were computed (Section~\ref{sec:deviations}). Cross-model attribution, window-length sensitivity and the 17-class task were listed as exploratory in the plan and were not run.
